\section*{5.4 AND/OR gadgets (author source pp.30--31)}
A function $g:\bits^m\to\bits$ affine projects to a function $h:\bits^p\to\bits$ if there are affine functions $\ell_1,\ldots,\ell_m:\mathbb Z_2^p\to\mathbb Z_2$ such that
\[g(\ell_1(z),\ldots,\ell_m(z))=h(z).\]
A gadget $g:\bits^m\to\bits$ is AND/OR if it affine projects to both binary AND and binary OR. Here are some examples of AND/OR gadgets:
\begin{itemize}
\item $g(x,y,z):=x\vee(y\wedge z)$. The projections: $g(0,a,b)=a\wedge b$ and $g(a,b,1)=a\vee b$.
\item $g(x,y,z):=[x+y+z=1]$. The projections: $g(1,\bar a,\bar b)=a\wedge b$ and $g(\bar a,\bar b,a\oplus b)=a\vee b$.
\end{itemize}
\begin{originalresult}{Theorem 5.12}
If $g$ is an AND/OR gadget then for all functions $f:\bits^n\to\bits$
\[\log\DTSize_\oplus(f\comp g)\ge C(f).\]
\end{originalresult}
\begin{proof}
Let $T_0$ be a parity decision tree for $f\comp g$. Let $z$ be a point such that $C(f,z)=C(f)$. Since $g$ is AND/OR, we can apply affine substitutions to $T_0$ and get a parity decision tree $T$ computing function $f\comp(h_{z_1},\ldots,h_{z_n})$, where $h_0(a,b):=a\wedge b$ and $h_1(a,b):=a\vee b$.
Now we follow the simulation proof of Chattopadhyay et al. (2023). We start at the root. When at a node $v$, if the query is constant then we follow the corresponding edge, otherwise we pick some variable, without loss of generality $a_i$, which appears in the query. So the query is of the form $a_i\oplus\ell$. We set $b_i:=0$ if $z_i=0$, and we set $b_i:=1$ if $z_i=1$. Then we choose $c$, the answer to the query, such that it leads to the smaller subtree, and set $a_i:=\ell\oplus c$. From now on, we see neither $a_i$ nor $b_i$.
Eventually, we reach a leaf. Suppose that this happens after eliminating coordinates $i_1,\ldots,i_K$, which means that $T$ contains at least $2^K$ leaves. For every input $z'$ to $f$ which agrees with $z$ on $i_1,\ldots,i_K$ we can find an input $w'$ to $T$ which satisfies $g(w')=z'$ and leads to the same leaf, showing that $z'_{i_1}=z_{i_1},\ldots,z'_{i_K}=z_{i_K}$ is a certificate for $f$ at $z$. We conclude that $K\ge C(f,z)=C(f)$, and so $|T_0|\ge|T|\ge2^K$.
\end{proof}
\section*{5.5 Gadget classification (author source pp.31--33)}
For the proof, we will need the following classical fact.
\begin{originalresult}{Proposition 5.13 (Chen2009, Theorem4.3)}
Suppose that $f:\bits^n\to\bits$ satisfies the following condition: whenever $f(x)=f(y)=f(z)=1$ then $f(x\oplus y\oplus z)=1$. Then either $f\equiv0$, or $f$ is an indicator of an affine subspace.
\end{originalresult}
We need this lemma to classify gadgets that are not AND/OR.
\begin{originalresult}{Lemma 5.14}
If $g$ is not an AND/OR then $g$ is a disjunction or a conjunction of affine forms in the inputs.
\end{originalresult}
\begin{proof}
Suppose that $g$ does not affine project to OR. We claim that for any inputs $x,y,z$, if $g(x)=g(y)=g(z)=1$ then $g(x\oplus y\oplus z)=1$. Indeed, suppose that $g(x)=g(y)=g(z)=1$ but $g(x\oplus y\oplus z)=0$, and consider the projection
\[h(a,b):=g(x\oplus y\oplus z\oplus a(x\oplus y)\oplus b(x\oplus z)).\]
Then $h(0,0)=g(x\oplus y\oplus z)=0$, $h(1,0)=g(z)=1$, $h(0,1)=g(y)=1$, and $h(1,1)=g(x)=1$. Thus, we can conclude that $g$ affine projects to OR.
Applying Proposition 5.13, we see that if $g$ does not affine project to OR then it is either the empty disjunction or a conjunction of affine forms. Similarly, if $g$ does not affine project to AND then $\bar g$ does not affine project to OR, and so $g$ is either the empty conjunction or a disjunction of affine forms.
\end{proof}
Now we can prove the classification of gadgets.
\begin{originalresult}{Theorem 5.15}
For every gadget $g$, one of the following cases holds:
\begin{enumerate}
\item[(i)] There is an infinite family of functions $f_n$ with $\DT(f_n)\to\infty$ and $\DTSize_\oplus(f_n\comp g)=O(1)$.
\item[(ii)] For every function $f$, $\DT_\oplus(f\comp g)=\Omega(\DT(f)^{\Omega(1)})$. There is an infinite family of functions $f_n$ with $\DT(f_n)\to\infty$ and $\DTSize(f_n\comp g)=O(\DT(f))$.
\item[(iii)] For every function $f$, $\log\DTSize_\oplus(f\comp g)=\Omega(\DT(f)^{\Omega(1)})$ and also $C_\oplus(f\comp g)=\Omega(C(f)^{\Omega(1)})$.
\end{enumerate}
\end{originalresult}
\begin{proof}
If gadget $g$ is AND/OR then by Theorem 5.12 we have
\[\log\DTSize_\oplus(f\comp g)\ge C(f)=\Omega(\sqrt{\DT(f)}).\]
Since $C_\oplus(f\comp g)=\Omega(\sqrt{\DT_\oplus(f\comp g)})$ (Zhang \& Shi2010) and also $\DT_\oplus(f\comp g)\ge\log\DTSize_\oplus(f\comp g)$, we get the third case of the theorem.
Otherwise, by Lemma 5.14, gadget $g$ is either a conjunction of affine forms or a disjunction of affine forms. Suppose, without loss of generality, that $g$ is a disjunction of affine forms. If $g\equiv0$ or $g$ is an affine form, then we can take $f_n=\oplus_n$ to get the first case of the theorem. If $g$ is not affine then it affine projects to OR, and thus we get the second case: we need to apply Theorem 5.7 together with the polynomial relations between (parity) certificate complexity and (parity) decision tree complexity; the infinite family of functions is $f_n=\OR_n$.
\end{proof}
